This section presents and analyses the results of our four experiments alongside
the official task baselines.

\subsection{Main Results}

Table~\ref{tab:main_results} presents the evaluation results for all four RADAR
configurations on the combined validation set and on the RAID test set.
Table~\ref{tab:pan26_results} reports performance on the PAN-2026 validation set
in isolation, allowing a direct comparison with the published baseline numbers.

\begin{table*}[t]
\caption{Results on the combined validation set (Val, 18,423 samples) and the
  RAID test set (RAID-Test, 10,000 samples). Best results per metric are \textbf{bolded}.}
\label{tab:main_results}
\begin{tabular}{llrrrrrr}
\toprule
\textbf{System} & \textbf{Split} & \textbf{ROC-AUC} & \textbf{Brier} & \textbf{C@1} & \textbf{F1} & \textbf{F0.5u} & \textbf{Mean} \\
\midrule
Exp.~1: Enc.-Freeze (PAN26+RAID)   & Val       & 0.976 & 0.940 & 0.919 & 0.921 & 0.946 & 0.941 \\
Exp.~1: Enc.-Freeze (PAN26+RAID)   & RAID-Test & 0.970 & 0.934 & 0.913 & 0.907 & 0.941 & 0.933 \\
\midrule
Exp.~2: End-to-End (PAN26+RAID)    & Val       & \textbf{0.992} & \textbf{0.973} & \textbf{0.971} & \textbf{0.973} & \textbf{0.981} & \textbf{0.978} \\
Exp.~2: End-to-End (PAN26+RAID)    & RAID-Test & \textbf{0.990} & \textbf{0.967} & \textbf{0.964} & \textbf{0.964} & \textbf{0.974} & \textbf{0.972} \\
\midrule
Exp.~3: Enc.-Freeze (PAN26 only)   & Val       & — & — & — & — & — & — \\
Exp.~3: Enc.-Freeze (PAN26 only)   & RAID-Test & 0.716 & 0.731 & 0.650 & 0.659 & 0.649 & 0.681 \\
\midrule
Exp.~4: Enc.-Freeze (RAID only)    & Val       & 0.972 & 0.937 & 0.914 & 0.913 & 0.929 & 0.933 \\
Exp.~4: Enc.-Freeze (RAID only)    & RAID-Test & 0.974 & 0.941 & 0.920 & 0.918 & 0.935 & 0.938 \\
\bottomrule
\end{tabular}
\end{table*}

\begin{table}[t]
\caption{Comparison with official task baselines on the PAN-2026 validation set
  (3,589 samples). Baselines are evaluated in-domain; RADAR results on this
  same split are shown for direct comparison.}
\label{tab:pan26_results}
\begin{tabular}{lrrrrrr}
\toprule
\textbf{System} & \textbf{ROC-AUC} & \textbf{Brier} & \textbf{C@1} & \textbf{F1} & \textbf{F0.5u} & \textbf{Mean} \\
\midrule
TF-IDF SVM          & 0.996 & 0.951 & 0.984 & 0.980 & 0.981 & 0.978 \\
PPMd CBC            & 0.786 & 0.799 & 0.757 & 0.812 & 0.778 & 0.786 \\
Binoculars          & 0.918 & 0.867 & 0.844 & 0.872 & 0.882 & 0.877 \\
\midrule
Exp.~1 (Enc.-Freeze) & 0.997 & 0.982 & 0.976 & 0.982 & 0.978 & 0.983 \\
Exp.~2 (End-to-End)  & \textbf{1.000} & \textbf{0.997} & \textbf{0.996} & \textbf{0.997} & \textbf{0.996} & \textbf{0.997} \\
Exp.~3 (PAN26 only)  & 0.997 & 0.983 & 0.979 & 0.984 & 0.982 & 0.985 \\
Exp.~4 (RAID only)   & 0.795 & 0.651 & 0.644 & 0.784 & 0.694 & 0.714 \\
\bottomrule
\end{tabular}
\end{table}

\subsection{Analysis}

\paragraph{Two-stage training yields the best overall performance.}
The end-to-end fine-tuned model (Exp.~2) achieves the strongest results across all metrics
and on all evaluation sets. On the PAN-2026 validation set, it reaches near-perfect performance
(Mean~=~0.997, ROC-AUC~=~1.000), substantially outperforming the TF-IDF SVM baseline
(Mean~=~0.978) despite the latter being a strong and well-calibrated supervised model.
On the harder combined validation set that includes adversarial RAID examples, Exp.~2
achieves Mean~=~0.978, matching the SVM's in-domain performance on a considerably more
difficult evaluation distribution.

\paragraph{Frozen encoder stage provides a strong intermediate.}
Experiment~1 (frozen encoder, combined data) achieves Mean~=~0.941 on the combined validation
set and Mean~=~0.933 on the RAID test set. These results confirm that the SLPE and fusion
modules already extract substantial discriminative signal from the frozen encoder
representations combined with stylometric features—even with only 1.2\% of parameters
being updated. The Stage~1 checkpoint provides a well-initialised starting point that
accelerates Stage~2 convergence.

\paragraph{Data composition critically affects cross-domain robustness.}
Experiment~3 (trained on PAN-2026 only) achieves excellent in-domain performance
(Mean~=~0.985 on PAN26-Val) but catastrophically low performance on the RAID test set
(Mean~=~0.681), dropping over 30 percentage points. This domain gap illustrates
the danger of fitting exclusively to the in-domain distribution when the evaluation
set contains new generators and adversarial obfuscations.

Conversely, Experiment~4 (trained on RAID only) is quite robust on the RAID test set
(Mean~=~0.938) but degrades significantly on the PAN-2026 in-domain validation
(Mean~=~0.714). This suggests that RAID alone is not representative enough to fully
cover the stylistic characteristics of PAN-2026 texts.

The full combination of both datasets (Exp.~1 and Exp.~2) achieves a balanced and
superior performance profile, confirming that multi-source heterogeneous training is
essential for genuine robustness.

\paragraph{RADAR outperforms zero-shot baselines by a substantial margin.}
Even the frozen-encoder Stage~1 model (Exp.~1) substantially outperforms both
Binoculars (Mean~=~0.877) and PPMd CBC (Mean~=~0.786) on the PAN-2026 validation set
(Exp.~1: Mean~=~0.983). This confirms that training a hybrid model that combines
semantic and stylometric signals is more effective than unsupervised statistical
approaches, even without full encoder fine-tuning.

\paragraph{Value of stylometric features.}
An implicit ablation is provided by contrasting Exp.~3 (PAN-2026 only, frozen encoder)
against the TF-IDF SVM. While both achieve comparable in-domain performance (Mean~=~0.983
vs.~0.978), RADAR's advantage emerges clearly under distribution shift. The TF-IDF SVM
is purely lexical and known to fail under synonym substitution and paraphrase; the RADAR
architecture's stylometric features are more stable under these perturbations because
they capture structural patterns rather than surface vocabulary.

\paragraph{Calibration.}
The Brier score, which measures probabilistic calibration, consistently improves from
Stage~1 to Stage~2. The end-to-end model achieves Brier~=~0.997 on PAN26-Val and
Brier~=~0.967 on the RAID test set, indicating that the confidence scores are
well-calibrated and reliable for downstream use—important for the ``undecidable''
abstention mechanism in the PAN evaluation.

\subsection{Training Dynamics}

Tables~\ref{tab:training_steps_exp1} and~\ref{tab:training_steps_exp2} present selected
checkpoint metrics for Experiments~1 and~2. In Stage~1
(frozen encoder), the training loss decreases from 0.398 to 0.038 over 10 epochs,
while the validation loss stabilises around 0.163, with the best checkpoint at epoch~9.5
(Mean~=~0.941). In Stage~2 (full fine-tuning), the training loss decreases substantially
faster from 0.210 to 0.003, reflecting the additional capacity of the unfrozen encoder,
while the validation loss reaches its minimum around epoch~9 before showing minor
overfitting signs. We select the best checkpoint by Mean score on the validation set.

\begin{table*}[t]
\caption{Selected training checkpoints for Experiment~1 (frozen encoder, combined data),
showing key validation metrics at representative epochs.}
\label{tab:training_steps_exp1}
\begin{tabular}{rrrrrrrr}
\toprule
\textbf{Step} & \textbf{Epoch} & \textbf{Train Loss} & \textbf{Val Loss} & \textbf{ROC-AUC} & \textbf{C@1} & \textbf{F1} & \textbf{Mean} \\
\midrule
500   & 0.41 & 0.398 & 0.326 & 0.900 & 0.816 & 0.806 & 0.854 \\
2500  & 2.08 & 0.242 & 0.217 & 0.961 & 0.885 & 0.883 & 0.915 \\
5000  & 4.15 & 0.090 & 0.186 & 0.971 & 0.903 & 0.902 & 0.930 \\
8000  & 6.64 & 0.079 & 0.164 & 0.976 & 0.919 & 0.922 & 0.940 \\
11500 & 9.55 & 0.051 & 0.164 & 0.976 & 0.918 & 0.920 & 0.940 \\
12000 & 9.97 & 0.038 & 0.163 & 0.976 & 0.918 & 0.920 & 0.940 \\
\bottomrule
\end{tabular}
\end{table*}

\begin{table*}[t]
\caption{Selected training checkpoints for Experiment~2 (end-to-end fine-tuning),
showing key validation metrics at representative epochs.}
\label{tab:training_steps_exp2}
\begin{tabular}{rrrrrrrr}
\toprule
\textbf{Step} & \textbf{Epoch} & \textbf{Train Loss} & \textbf{Val Loss} & \textbf{ROC-AUC} & \textbf{C@1} & \textbf{F1} & \textbf{Mean} \\
\midrule
500   & 0.41 & 0.210 & 0.181 & 0.982 & 0.904 & 0.902 & 0.934 \\
2000  & 1.66 & 0.048 & 0.087 & 0.993 & 0.961 & 0.962 & 0.973 \\
4000  & 3.32 & 0.035 & 0.136 & 0.989 & 0.955 & 0.957 & 0.967 \\
8500  & 7.06 & 0.013 & 0.146 & 0.991 & 0.962 & 0.964 & 0.972 \\
11000 & 9.14 & 0.011 & 0.135 & 0.992 & 0.971 & 0.973 & 0.978 \\
12000 & 9.97 & 0.003 & 0.135 & 0.992 & 0.971 & 0.973 & 0.978 \\
\bottomrule
\end{tabular}
\end{table*}
